% =====================================================================
\chapter{Bối cảnh nghiệp vụ và yêu cầu}\label{chap:req}
% =====================================================================

\section{Bài toán}
Hệ thống quan sát \textbf{trạng thái công khai của catalog và giá} trên sàn thương mại điện tử Việt Nam.
Tiki là nguồn chính; Lazada là nguồn minh hoạ khả năng đa nguồn. Dữ liệu có được là những gì sàn hiển thị
cho một khách vãng lai trên trang danh mục: tiêu đề, giá, giá niêm yết, mức giảm, tình trạng còn hàng,
điểm đánh giá, các bộ đếm công khai (đã bán, số đánh giá) và vị trí xếp hạng.

Hệ thống \textbf{không} có đơn hàng, thanh toán, giỏ hàng hay người dùng, và không mô hình nào giả lập chúng.
Bộ đếm \code{sold_count} tăng không có nghĩa là một đơn hàng; mọi bảng chạm tới counter đều giữ cả delta thô
lẫn phán quyết hợp lệ của từng bước chuyển.

\section{Câu hỏi nghiệp vụ}
\begin{xltabular}{\linewidth}{@{}L{8mm}Y L{58mm}@{}}
  \caption{Câu hỏi nghiệp vụ và nơi trả lời}\label{tab:bq}\\
  \toprule
  \textbf{\#} & \textbf{Câu hỏi} & \textbf{Trả lời ở} \\
  \midrule
  BQ1 & Giá của một offer thay đổi thế nào theo thời gian? Thay đổi nào bất thường so với lịch sử của chính nó? &
        mart \code{offer_price_history_daily}, \code{price_anomaly_daily} \\
  BQ2 & Mức giảm giá sàn hiển thị có đáng tin không, so với giá trước đó của chính offer? &
        mart \code{discount_credibility_daily}, \code{discount_daily} (ADR-0009) \\
  BQ3 & Khi một listing thay đổi (giá, tồn kho, đánh giá), bao lâu thì thấy được? Độ phủ và độ tươi của mỗi nguồn ra sao? &
        speed layer; mart \code{source_coverage_daily}, \code{offer_freshness}, \code{crawl_reliability_daily} \\
  \bottomrule
\end{xltabular}

\section{Bên liên quan và người dùng}
\begin{xltabular}{\linewidth}{@{}L{38mm}Y L{40mm}@{}}
  \toprule
  \textbf{Vai trò} & \textbf{Nhu cầu} & \textbf{Giao diện} \\
  \midrule
  Nhà phân tích giá & lịch sử giá, giá bất thường, độ tin cậy khuyến mãi theo offer và danh mục & Superset (dashboard \code{marketplace-batch-gold}) \\
  Người theo dõi thời gian thực & thay đổi giá, tồn kho, giảm giá lớn ngay khi crawler thấy & Kibana, Redis \code{rt:*}, topic \code{marketplace.changes.v1} \\
  Người vận hành & crawl có chạy không, batch có promote không, quality có chặn không & Grafana, Airflow UI, \code{scripts/mp.ps1 status} \\
  Kỹ sư dữ liệu & thêm nguồn, sửa normalizer, chạy lại lịch sử & CLI \code{ecomdata}, DAG \code{marketplace_reprocess} \\
  \bottomrule
\end{xltabular}

\section{Yêu cầu chức năng}
\begin{xltabular}{\linewidth}{@{}L{12mm}Y L{44mm}@{}}
  \caption{Yêu cầu chức năng}\label{tab:fr}\\
  \toprule
  \textbf{Mã} & \textbf{Yêu cầu} & \textbf{Hiện thực} \\
  \midrule
  \endfirsthead
  \toprule \textbf{Mã} & \textbf{Yêu cầu} & \textbf{Hiện thực} \\ \midrule
  \endhead
  FR-01 & Thu thập trang danh mục theo lịch, có ưu tiên (tier), tôn trọng robots.txt và giới hạn tốc độ theo host. & frontier, \code{infra/http.py} (Chương~\ref{chap:ingest}) \\
  FR-02 & Lưu nguyên byte mọi phản hồi, kể cả phản hồi lỗi, kèm checksum. & Bronze/raw \\
  FR-03 & Giữ mọi trường của nguồn (trừ PII khai báo), không chọn trường lúc crawl (ELT). & Bronze/records, ADR-0001 \\
  FR-04 & Chuẩn hoá mọi nguồn về một mô hình observation chung; tách dòng lỗi vào quarantine có lý do. & stage \code{silver} \\
  FR-05 & Dựng star schema có SCD2 cho listing, dùng chung cho mọi run. & stage \code{dimensions}, \code{facts} \\
  FR-06 & Dựng 12 mart theo run, cắt tại \code{as_of}, chỉ xét dòng organic. & stage \code{marts} \\
  FR-07 & Chặn publish khi một luật chất lượng bắt buộc không đạt; phục vụ bản tốt cuối cùng. & stage \code{quality}, \code{publish}, \code{promote} \\
  FR-08 & Phát hiện 7 loại thay đổi theo offer trong thời gian gần thực. & speed layer (Mục~\ref{sec:speed}) \\
  FR-09 & Truy ngược từ một con số trên dashboard về byte HTTP gốc. & lineage (Mục~\ref{sec:lineage}) \\
  FR-10 & Chạy lại một khoảng lịch sử hoặc dựng lại toàn bộ từ Bronze khi parser/normalizer đổi. & \code{reparse}, DAG \code{marketplace_reprocess} \\
  FR-11 & Thêm một nguồn mới mà không sửa batch, speed hay ops. & \code{SourceSpec}, ADR-0004 \\
  FR-12 & Sao lưu và khôi phục nhất quán giữa lake và cache. & \code{ops backup}/\code{restore} \\
  \bottomrule
\end{xltabular}

\section{Yêu cầu phi chức năng}
\begin{xltabular}{\linewidth}{@{}L{14mm}L{28mm}Y L{36mm}@{}}
  \caption{Yêu cầu phi chức năng và mức đo được}\label{tab:nfr}\\
  \toprule
  \textbf{Mã} & \textbf{Thuộc tính} & \textbf{Mục tiêu} & \textbf{Đo được} \\
  \midrule
  \endfirsthead
  \toprule \textbf{Mã} & \textbf{Thuộc tính} & \textbf{Mục tiêu} & \textbf{Đo được} \\ \midrule
  \endhead
  NFR-01 & Độ tươi batch & mart được promote mỗi giờ; cảnh báo khi run promote mới nhất cũ hơn 2,5 giờ & 10/10 run PROMOTED trong 9 giờ 16 phút (13--14/11) \\
  NFR-02 & Độ trễ speed & thay đổi xuất hiện trên Kibana trong khoảng 1 phút sau khi crawl (trigger 30 s) & p95 khoảng 59--62 s sau khi đuổi kịp tồn đọng \\
  NFR-03 & Độ đúng & batch và speed cùng định nghĩa observation; chạy lại không sinh bản trùng & một normalizer dùng chung; ID tất định \\
  NFR-04 & Khả năng phục hồi & mất Kafka, ES, Redis hoặc PostgreSQL không mất dữ liệu gốc & Bronze là nguồn sự thật; ma trận lỗi Mục~\ref{sec:failure-matrix} \\
  NFR-05 & Tái lập & một run tái lập được trên đúng star mà nó đã đọc & manifest ghi \code{dw_versions} \\
  NFR-06 & Lịch sự với site & $\ge$ 2 s/request/host (Tiki), $\ge$ 5 s (Lazada); không vượt anti-bot & \code{Throttle}, robots, circuit breaker \\
  NFR-07 & Quan sát được & mọi tầng có metrics, log, cảnh báo & 5 dashboard, 19 luật cảnh báo \\
  NFR-08 & Tài nguyên & chạy được trên một laptop, mỗi lượt $\le$ 5--6 giờ & stack local một bản mỗi service \\
  NFR-09 & Khả chuyển & cùng mã chạy được trên cloud & job Databricks (DBR 17.3) \\
  NFR-10 & Bảo trì được & luật import giữa các tầng được kiểm trong CI & \code{import-linter}, \code{tests/architecture/} \\
  \bottomrule
\end{xltabular}

\section{Ràng buộc và giả định}
\begin{itemize}
  \item \textbf{Pháp lý, đạo đức thu thập}: chỉ dữ liệu công khai; tôn trọng robots.txt (không tải được robots thì coi như cấm toàn bộ);
        user-agent tự khai danh tính (\code{EcommercePriceResearchBot/0.1}); không vượt captcha hay chặn traffic.
        Vì vậy Shopee không được crawl (ADR-0006).
  \item \textbf{Phần cứng}: stack chạy trên một laptop/PC cá nhân; Docker Desktop; mỗi lượt chạy liên tục tối đa khoảng 5--6 giờ.
  \item \textbf{Phạm vi dữ liệu}: Tiki chỉ trang 1 của mỗi danh mục (40 sản phẩm); Lazada có phân trang nhưng bị chặn nhanh.
  \item \textbf{Lưu trữ}: lake (Delta trên MinIO/S3) là bản lưu; PostgreSQL chỉ là cache phục vụ và kho trạng thái vận hành (quyết định 26/10/2026, ADR-0008).
  \item \textbf{Thời gian}: mọi mốc thời gian dữ liệu là UTC; \code{run_id} và \code{as_of} theo UTC.
\end{itemize}

\section{Ngoài phạm vi}
\begin{itemize}
  \item Đơn hàng, doanh thu thật, hành vi người dùng.
  \item Dự báo nhu cầu theo danh mục: hoãn tới khi có ít nhất ba tuần dữ liệu liên tục (ADR-0009).
  \item Ghép cùng một sản phẩm giữa các sàn (product matching).
  \item Chạy 24/7 có tính sẵn sàng cao: mọi service chạy một bản.
\end{itemize}
